\chapter{Descriptive complexity}
\label{chap:descriptive}

Descriptive complexity characterizes complexity classes by the logics that
define their problems, after Immerman and Fagin. Once a class is captured by a
logic, a lower bound becomes a question of what the logic can express. The
development lives in \texttt{Complexitylib.DescriptiveComplexity}. The following
are formalized:
\begin{itemize}
  \item finite structures over finite vocabularies, their isomorphisms, and
    injective homomorphisms and substructure inclusions between them;
  \item first-order logic with de Bruijn variables, substitution, and invariance
    under isomorphism;
  \item Boolean queries, order independence, and first-order definability;
  \item second-order logic, its $\exists\cc{SO}$ fragment, and invariance under
    isomorphism;
  \item a bit-string encoding of structures and the language a query induces;
  \item a computable first-order model checker, proved correct, with no bound on
    its running time;
  \item universe-preserving (dimension-1) first-order reductions and their
    quantifier-free restriction, both of which are preorders.
\end{itemize}
No capture theorem is formalized yet, and the link from logic to machines stops
at the definition of the language of a query: there is no decoder, no running-time
bound, and no machine construction. The main planned result is Fagin's theorem
$\cc{NP} = \exists\cc{SO}$. Later planned results are $\cc{FO} \subseteq \cc{AC^0}$
and the Immerman--Vardi theorem.

\paragraph{Conventions.} A structure has universe $\mathrm{Fin}(n)$ with
$n \ge 2$ (Immerman's Proviso~1.15). Queries are properties of structures with
$\mathrm{Prop}$-valued relations. The encoding and the model checker take
decidable structures, whose relations are $\mathrm{Bool}$-valued and which coerce
to ordinary structures. The canonical order on $\mathrm{Fin}(n)$ is available as
meta-level helpers, but the first- and second-order syntax is unordered, so
formulas cannot mention it. Order independence is therefore a theorem about the
existing logics rather than a restriction imposed on them.

\section{Vocabularies and finite structures}

\begin{definition}[Vocabulary]
  \label{def:desc-vocabulary}
  \lean{Complexity.DescriptiveComplexity.Vocabulary, Complexity.DescriptiveComplexity.Vocabulary.IsRelational, Complexity.DescriptiveComplexity.Vocabulary.graph}
  \leanok
  A vocabulary has $r$ relation symbols, indexed by $\mathrm{Fin}(r)$, each with
  an arity in $\mathbb{N}$, and $c$ constant symbols, indexed by
  $\mathrm{Fin}(c)$. It is relational if $c = 0$. The graph vocabulary has one
  binary relation symbol and no constants.
\end{definition}

\begin{definition}[Finite structure]
  \label{def:desc-structure}
  \lean{Complexity.DescriptiveComplexity.FinStruct, Complexity.DescriptiveComplexity.DecFinStruct, Complexity.DescriptiveComplexity.DecFinStruct.toFinStruct, Complexity.DescriptiveComplexity.FinStruct.leRel, Complexity.DescriptiveComplexity.FinStruct.sucRel, Complexity.DescriptiveComplexity.FinStruct.minElem, Complexity.DescriptiveComplexity.FinStruct.maxElem}
  \leanok
  \uses{def:desc-vocabulary}
  A finite structure $\mathcal{A}$ over $V$ consists of:
  \begin{itemize}
    \item a size $n \ge 2$ and the universe $\mathrm{Fin}(n)$;
    \item for each relation symbol of arity $k$, a $\mathrm{Prop}$-valued
      relation on $k$-tuples, that is, on maps
      $\mathrm{Fin}(k) \to \mathrm{Fin}(n)$;
    \item for each constant symbol, an element.
  \end{itemize}
  A decidable structure has the same data, with $\mathrm{Bool}$-valued
  relations, and coerces to a structure in which a relation holds where its value
  is $\mathrm{true}$. The canonical order $\le$, the successor relation
  ($b = a + 1$), the minimum $0$, and the maximum $n - 1$ of $\mathrm{Fin}(n)$
  are defined for every structure as meta-level helpers. They are not part of the
  structure, and no formula can mention them.
\end{definition}

\begin{definition}[Isomorphisms, embeddings, injective homomorphisms]
  \label{def:desc-iso}
  \lean{Complexity.DescriptiveComplexity.Iso, Complexity.DescriptiveComplexity.Iso.refl,
    Complexity.DescriptiveComplexity.Iso.symm, Complexity.DescriptiveComplexity.Iso.trans,
    Complexity.DescriptiveComplexity.Embedding, Complexity.DescriptiveComplexity.Embedding.refl,
    Complexity.DescriptiveComplexity.Embedding.trans, Complexity.DescriptiveComplexity.Embedding.ofIso,
    Complexity.DescriptiveComplexity.Embedding.toInjectiveHom,
    Complexity.DescriptiveComplexity.InjectiveHom, Complexity.DescriptiveComplexity.InjectiveHom.ofIso}
  \leanok
  \uses{def:desc-structure}
  An isomorphism $\mathcal{A} \cong \mathcal{B}$ is a pair of mutually inverse
  maps between the universes whose forward map preserves and reflects every
  relation and preserves constants. Isomorphisms have identities, inverses, and
  composites. An \emph{embedding} is an injective map that preserves and
  reflects every relation and preserves constants, as in model theory; its
  image is a substructure of $\mathcal{B}$ isomorphic to $\mathcal{A}$.
  Embeddings have identities and composites, and every isomorphism is one. An
  \emph{injective homomorphism} only preserves relations, so $\mathcal{B}$ may
  relate images of unrelated elements; every embedding is one. All three
  notions carry their maps as data.
\end{definition}

\begin{theorem}[Isomorphic structures have equal size]
  \label{thm:desc-iso-card}
  \lean{Complexity.DescriptiveComplexity.Iso.card_eq}
  \leanok
  \uses{def:desc-iso}
  If $\mathcal{A} \cong \mathcal{B}$, then $|\mathcal{A}| = |\mathcal{B}|$.
\end{theorem}
\begin{proof}
  \leanok
  Both maps are injections between $\mathrm{Fin}$ types, so the pigeonhole
  principle applies in each direction.
\end{proof}

\section{First-order logic}

\begin{definition}[First-order syntax]
  \label{def:desc-fo-syntax}
  \lean{Complexity.DescriptiveComplexity.Term, Complexity.DescriptiveComplexity.Formula, Complexity.DescriptiveComplexity.Sentence, Complexity.DescriptiveComplexity.Formula.quantifierRank, Complexity.DescriptiveComplexity.Formula.size}
  \leanok
  \uses{def:desc-vocabulary}
  Terms with $k$ free variables are de Bruijn variables $x_i$ ($i < k$) or
  constant symbols. Formulas are built from relation atoms, equality, $\neg$,
  $\wedge$, $\vee$, $\exists$, and $\forall$. A quantifier binds variable $0$ and
  shifts the others. Sentences are formulas with no free variables. Quantifier
  rank (the quantifier nesting depth) and size (the number of nodes) are defined
  by structural recursion.
\end{definition}

\begin{definition}[First-order semantics]
  \label{def:desc-fo-semantics}
  \lean{Complexity.DescriptiveComplexity.Env, Complexity.DescriptiveComplexity.envCons, Complexity.DescriptiveComplexity.emptyEnv, Complexity.DescriptiveComplexity.Term.eval, Complexity.DescriptiveComplexity.Formula.Sat, Complexity.DescriptiveComplexity.Sentence.Models}
  \leanok
  \uses{def:desc-fo-syntax, def:desc-structure}
  An environment for $k$ free variables is a map
  $\mathrm{Fin}(k) \to \mathrm{Fin}(|\mathcal{A}|)$. Extending it by an element
  puts that element at index $0$ and shifts the others. Satisfaction
  $\mathcal{A}, \sigma \models \varphi$ is Tarskian, and quantifiers range over
  $\mathrm{Fin}(|\mathcal{A}|)$. $\mathcal{A} \models \varphi$ means
  satisfaction of a sentence under the empty environment.
\end{definition}

\begin{theorem}[Isomorphism invariance of first-order logic]
  \label{thm:desc-fo-iso-invariance}
  \lean{Complexity.DescriptiveComplexity.Term.eval_iso, Complexity.DescriptiveComplexity.Formula.sat_iso}
  \leanok
  \uses{def:desc-fo-semantics, def:desc-iso}
  Let $f : \mathcal{A} \cong \mathcal{B}$. For every term $t$ and environment
  $\sigma$, the value of $t$ in $\mathcal{B}$ under $f \circ \sigma$ is the image
  under $f$ of its value in $\mathcal{A}$ under $\sigma$. For every formula
  $\varphi$ and environment $\sigma$, $\mathcal{A}, \sigma \models \varphi$ iff
  $\mathcal{B}, f \circ \sigma \models \varphi$.
\end{theorem}
\begin{proof}
  \leanok
  Terms are variables or constants, and $f$ preserves constants. For formulas,
  induct on $\varphi$. Quantifier cases move witnesses along $f$ and $f^{-1}$.
\end{proof}

\begin{definition}[Substitution]
  \label{def:desc-substitution}
  \lean{Complexity.DescriptiveComplexity.Term.shift, Complexity.DescriptiveComplexity.Term.substTerm, Complexity.DescriptiveComplexity.liftSubst, Complexity.DescriptiveComplexity.Formula.subst}
  \leanok
  \uses{def:desc-fo-syntax}
  Simultaneous substitution $\varphi[\rho]$ takes a formula with $m$ free
  variables and terms $\rho(0), \dots, \rho(m-1)$ with $n$ free variables, and
  replaces each $x_k$ by $\rho(k)$, giving a formula with $n$ free variables.
  Under a binder, the substitution is lifted: variable $0$ stays, and every other
  image is shifted up one de Bruijn index.
\end{definition}

\begin{theorem}[Substitution lemma]
  \label{thm:desc-substitution}
  \lean{Complexity.DescriptiveComplexity.Formula.subst_sat, Complexity.DescriptiveComplexity.Formula.quantifierRank_subst}
  \leanok
  \uses{def:desc-substitution, def:desc-fo-semantics}
  For every structure $\mathcal{A}$, environment $\sigma$, substitution $\rho$,
  and formula $\varphi$,
  $\mathcal{A}, \sigma \models \varphi[\rho]$ iff
  $\mathcal{A}, (k \mapsto \rho(k)^{\mathcal{A},\sigma}) \models \varphi$.
  Substitution preserves quantifier rank exactly.
\end{theorem}
\begin{proof}
  \leanok
  Induct on $\varphi$, generalizing over the number of free variables, and use a
  lemma evaluating lifted substitutions.
\end{proof}

\section{Boolean queries and definability}

\begin{definition}[Boolean query]
  \label{def:desc-query}
  \lean{Complexity.DescriptiveComplexity.BooleanQuery, Complexity.DescriptiveComplexity.BooleanQuery.IsOrderIndependent, Complexity.DescriptiveComplexity.BooleanQuery.complement, Complexity.DescriptiveComplexity.BooleanQuery.inter, Complexity.DescriptiveComplexity.BooleanQuery.union}
  \leanok
  \uses{def:desc-structure, def:desc-iso}
  A Boolean query over $V$ is a property of finite $V$-structures. It is
  \emph{order-independent} if isomorphic structures agree on it, that is, if it
  is invariant under isomorphism. Complement, intersection, and union are defined
  pointwise.
\end{definition}

\begin{proposition}[Closure of order independence]
  \label{prop:desc-order-independent-closure}
  \lean{Complexity.DescriptiveComplexity.BooleanQuery.complement_orderIndependent, Complexity.DescriptiveComplexity.BooleanQuery.inter_orderIndependent, Complexity.DescriptiveComplexity.BooleanQuery.union_orderIndependent}
  \leanok
  \uses{def:desc-query}
  Order-independent queries are closed under complement, intersection, and union.
\end{proposition}
\begin{proof}
  \leanok
  Immediate from the definitions.
\end{proof}

\begin{theorem}[First-order sentences are order-independent]
  \label{thm:desc-fo-order-independent}
  \lean{Complexity.DescriptiveComplexity.Sentence.orderIndependent}
  \leanok
  \uses{def:desc-query, def:desc-fo-semantics}
  For every sentence $\varphi$, the query $\mathcal{A} \mapsto (\mathcal{A} \models \varphi)$
  is order-independent (Immerman, Proposition~1.16).
\end{theorem}
\begin{proof}
  \leanok
  \uses{thm:desc-fo-iso-invariance}
  Specialize isomorphism invariance to the empty environment. As a worked check,
  \texttt{cycle3\_fo\_indistinguishable} applies this to two explicitly
  isomorphic directed 3-cycles, which therefore satisfy the same sentences.
\end{proof}

\begin{definition}[First-order definable query]
  \label{def:desc-fo-definable}
  \lean{Complexity.DescriptiveComplexity.FODefinable}
  \leanok
  \uses{def:desc-query, def:desc-fo-semantics}
  $Q$ is first-order definable if some sentence $\varphi$ satisfies
  $Q(\mathcal{A}) \iff \mathcal{A} \models \varphi$ for every $\mathcal{A}$.
\end{definition}

\begin{theorem}[First-order definable queries]
  \label{thm:desc-fo-definable}
  \lean{Complexity.DescriptiveComplexity.FODefinable.orderIndependent, Complexity.DescriptiveComplexity.FODefinable.complement, Complexity.DescriptiveComplexity.FODefinable.inter, Complexity.DescriptiveComplexity.FODefinable.union}
  \leanok
  \uses{def:desc-fo-definable}
  Every first-order definable query is order-independent. First-order definable
  queries are closed under complement, intersection, and union.
\end{theorem}
\begin{proof}
  \leanok
  \uses{thm:desc-fo-order-independent}
  Use the defining sentence, and $\neg$, $\wedge$, $\vee$ for the closure
  properties.
\end{proof}

\section{Second-order logic}

\begin{definition}[Second-order syntax]
  \label{def:desc-so-syntax}
  \lean{Complexity.DescriptiveComplexity.SOFormula, Complexity.DescriptiveComplexity.SOSentence, Complexity.DescriptiveComplexity.SOFormula.ofFormula, Complexity.DescriptiveComplexity.SOFormula.size}
  \leanok
  \uses{def:desc-fo-syntax}
  Second-order formulas extend first-order formulas with two constructs:
  applications of relation variables, and second-order quantifiers $\exists R$,
  $\forall R$ over a fresh relation of a given arity. Relation variables are
  de Bruijn-indexed by a context listing their arities, innermost first. A
  second-order sentence has an empty relation context and no free element
  variables. First-order formulas embed into second-order formulas over any
  relation context. Size counts nodes.
\end{definition}

\begin{definition}[Second-order semantics]
  \label{def:desc-so-semantics}
  \lean{Complexity.DescriptiveComplexity.REnv, Complexity.DescriptiveComplexity.rCons, Complexity.DescriptiveComplexity.emptyREnv, Complexity.DescriptiveComplexity.SOFormula.Sat, Complexity.DescriptiveComplexity.SOSentence.Models}
  \leanok
  \uses{def:desc-so-syntax, def:desc-fo-semantics}
  A relation environment assigns to each relation variable in the context a
  $\mathrm{Prop}$-valued relation of its arity on the universe. Extending it puts
  the new relation at index $0$. Satisfaction is defined under an element
  environment and a relation environment, and second-order quantifiers range over
  all $\mathrm{Prop}$-valued relations of the quantified arity on the universe.
  $\mathcal{A} \models \varphi$ for a second-order sentence means satisfaction
  under the empty environments.
\end{definition}

\begin{definition}[The $\exists\cc{SO}$ fragment]
  \label{def:desc-existential-so}
  \lean{Complexity.DescriptiveComplexity.SOFormula.IsFOMatrix, Complexity.DescriptiveComplexity.SOFormula.IsExistSO}
  \leanok
  \uses{def:desc-so-syntax}
  A formula is an \emph{FO matrix} if it has no second-order quantifiers, though
  it may apply relation variables and use first-order quantifiers. It is in
  $\exists\cc{SO}$ form if it is a possibly empty block of second-order
  existential quantifiers followed by an FO matrix.
\end{definition}

\begin{definition}[Second-order definable query]
  \label{def:desc-so-definable}
  \lean{Complexity.DescriptiveComplexity.SODefinable}
  \leanok
  \uses{def:desc-query, def:desc-so-semantics}
  $Q$ is second-order definable if some second-order sentence $\varphi$ satisfies
  $Q(\mathcal{A}) \iff \mathcal{A} \models \varphi$ for every $\mathcal{A}$.
\end{definition}

\begin{theorem}[First-order logic inside second-order logic]
  \label{thm:desc-fo-in-so}
  \lean{Complexity.DescriptiveComplexity.SOFormula.ofFormula_sat, Complexity.DescriptiveComplexity.SOSentence.models_ofFormula, Complexity.DescriptiveComplexity.SOFormula.ofFormula_isFOMatrix, Complexity.DescriptiveComplexity.SOFormula.ofFormula_isExistSO, Complexity.DescriptiveComplexity.FODefinable.toSODefinable}
  \leanok
  \uses{def:desc-so-semantics, def:desc-existential-so, def:desc-so-definable, def:desc-fo-definable}
  For every relation environment $\rho$, the embedding of a first-order formula
  holds at $(\sigma, \rho)$ exactly when the original holds at $\sigma$. In
  particular, a structure models the embedding of a sentence iff it models the
  sentence. The embedding is an FO matrix, hence in $\exists\cc{SO}$ form with an
  empty quantifier block. Consequently every first-order definable query is
  second-order definable.
\end{theorem}
\begin{proof}
  \leanok
  Induct on the first-order formula.
\end{proof}

\begin{theorem}[Isomorphism invariance of second-order logic]
  \label{thm:desc-so-iso-invariance}
  \lean{Complexity.DescriptiveComplexity.REnv.map, Complexity.DescriptiveComplexity.SOFormula.sat_iso, Complexity.DescriptiveComplexity.SOSentence.models_iso}
  \leanok
  \uses{def:desc-so-semantics, def:desc-iso}
  Along $f : \mathcal{A} \cong \mathcal{B}$, a second-order formula holds at
  $(\sigma, \rho)$ iff it holds at $(f \circ \sigma, f_*\rho)$, where
  $f_*\rho$ transports each relation along $f$:
  $(f_*\rho)(R)(\bar b) = \rho(R)(f^{-1} \circ \bar b)$. In particular,
  isomorphic structures satisfy the same second-order sentences.
\end{theorem}
\begin{proof}
  \leanok
  \uses{thm:desc-fo-iso-invariance}
  Induct on the formula. For second-order quantifiers, push relations forward
  along $f$ and pull them back along $f^{-1}$.
\end{proof}

\begin{theorem}[Second-order definable queries]
  \label{thm:desc-so-definable}
  \lean{Complexity.DescriptiveComplexity.SODefinable.orderIndependent, Complexity.DescriptiveComplexity.SODefinable.complement, Complexity.DescriptiveComplexity.SODefinable.inter, Complexity.DescriptiveComplexity.SODefinable.union}
  \leanok
  \uses{def:desc-so-definable}
  Every second-order definable query is order-independent. In particular, so is
  every query defined by a sentence in $\exists\cc{SO}$ form. Second-order
  definable queries are closed under complement, intersection, and union.
\end{theorem}
\begin{proof}
  \leanok
  \uses{thm:desc-so-iso-invariance}
  Specialize to the empty environments, and use $\neg$, $\wedge$, $\vee$ for the
  closure properties.
\end{proof}

\section{Encodings and model checking}

\begin{definition}[Encoding structures as bit strings]
  \label{def:desc-encoding}
  \lean{Complexity.DescriptiveComplexity.allTuples, Complexity.DescriptiveComplexity.encodeRelC, Complexity.DescriptiveComplexity.encodeRelsC, Complexity.DescriptiveComplexity.encodeConstC, Complexity.DescriptiveComplexity.encodeConstsC, Complexity.DescriptiveComplexity.encodeStruct}
  \leanok
  \uses{def:desc-structure}
  A decidable structure of size $n$ is encoded, computably, as the concatenation
  of three blocks:
  \begin{itemize}
    \item the size in unary, $1^n 0$;
    \item for each relation symbol in index order, its truth table, listed over a
      computable enumeration of all tuples of its arity;
    \item for each constant symbol in index order, a one-hot block of length $n$
      whose bit $i$ is $1$ exactly when the constant is $i$.
  \end{itemize}
  The encoding depends on the canonical order of $\mathrm{Fin}(n)$.
\end{definition}

\begin{proposition}[Encoding lengths]
  \label{prop:desc-encoding-length}
  \lean{Complexity.DescriptiveComplexity.allTuples_length, Complexity.DescriptiveComplexity.mem_allTuples, Complexity.DescriptiveComplexity.encodeRelC_length, Complexity.DescriptiveComplexity.encodeRelsC_length, Complexity.DescriptiveComplexity.encodeConstC_length, Complexity.DescriptiveComplexity.encodeConstsC_length, Complexity.DescriptiveComplexity.encodeStruct_length, Complexity.DescriptiveComplexity.encodeStruct_of_isRelational, Complexity.DescriptiveComplexity.encodeStruct_card, Complexity.DescriptiveComplexity.encodeConstC_injective}
  \leanok
  \uses{def:desc-encoding}
  The enumeration of $k$-tuples has length $n^k$ and contains every $k$-tuple,
  so a truth table of arity $k$ has length $n^k$. A one-hot block has length
  $n$. The encoding has length
  $n + 1 + \sum_R n^{\mathrm{ar}(R)} + c \cdot n$, where $c$ is the number of
  constants, and for a relational vocabulary it is the unary block followed by
  the truth tables alone. The size $n$ is recoverable as the length of the
  leading block of ones, and a one-hot constant block determines its element.
\end{proposition}
\begin{proof}
  \leanok
  Induct on the arity, then sum block lengths.
\end{proof}

\begin{definition}[Language of a query]
  \label{def:desc-query-language}
  \lean{Complexity.DescriptiveComplexity.queryLanguage, Complexity.DescriptiveComplexity.mem_queryLanguage}
  \leanok
  \uses{def:desc-query, def:desc-encoding, def:language}
  The language of $Q$ is
  $\{\mathrm{enc}(\mathcal{A}) : \mathcal{A} \text{ decidable}, Q(\mathcal{A})\}$,
  where $Q$ is applied to the structure that $\mathcal{A}$ coerces to. By
  definition every string in it is an encoding, and the encoding of a decidable
  structure satisfying $Q$ lies in it. This is the bridge from queries to the
  machine-model languages.
\end{definition}

\begin{definition}[Structure decoder]
  \label{def:desc-decoder}
  \uses{def:desc-encoding}
  A computable parser $\mathrm{dec} : \{0,1\}^* \to V\text{-structures} \cup \{\bot\}$.
  It reads the unary size and then the fixed-length relation and constant blocks,
  and returns $\bot$ on malformed input, including a size below $2$ and a
  constant block that is not one-hot.
\end{definition}

\begin{theorem}[Encoding round trip]
  \label{thm:desc-encoding-roundtrip}
  \uses{def:desc-decoder, def:desc-encoding}
  \notready
  $\mathrm{dec}(\mathrm{enc}(\mathcal{A})) = \mathcal{A}$, up to extensional
  equality of relations. Hence $\mathrm{enc}$ is injective, and
  $x \in \mathrm{queryLanguage}(Q)$ iff $\mathrm{dec}(x) = \mathcal{A} \neq \bot$
  and $Q(\mathcal{A})$.
\end{theorem}
\begin{proof}
  \uses{prop:desc-encoding-length}
  Recover the size with \texttt{encodeStruct\_card}, then split the remainder by
  the known block lengths.
\end{proof}

\begin{definition}[Boolean model checker]
  \label{def:desc-model-checking}
  \lean{Complexity.DescriptiveComplexity.Formula.evalB, Complexity.DescriptiveComplexity.Sentence.evalB}
  \leanok
  \uses{def:desc-fo-syntax, def:desc-structure}
  A computable $\mathrm{Bool}$-valued evaluator for first-order formulas on
  decidable structures. Atoms look up the $\mathrm{Bool}$-valued relations, and
  quantifiers become \texttt{any} and \texttt{all} over the list of universe
  elements. The sentence form runs it under the empty environment.
\end{definition}

\begin{theorem}[Model checking is correct]
  \label{thm:desc-model-checking}
  \lean{Complexity.DescriptiveComplexity.Formula.evalB_eq_sat, Complexity.DescriptiveComplexity.Sentence.evalB_eq_models, Complexity.DescriptiveComplexity.decidableModels}
  \leanok
  \uses{def:desc-model-checking, def:desc-fo-semantics}
  For every decidable structure $\mathcal{A}$, formula $\varphi$, and
  environment $\sigma$, $\mathrm{evalB}(\mathcal{A}, \sigma, \varphi) = \mathrm{true}$
  iff $\mathcal{A}, \sigma \models \varphi$, and likewise for sentences under the
  empty environment. Hence $\mathcal{A} \models \varphi$ is decidable for every
  decidable structure $\mathcal{A}$ and sentence $\varphi$. This is correctness
  only: no bound on the evaluator's running time and no machine implementation
  are formalized.
\end{theorem}
\begin{proof}
  \leanok
  Induct on $\varphi$.
\end{proof}

\begin{theorem}[First-order queries are in polynomial time]
  \label{thm:desc-fo-in-P}
  \uses{def:desc-query-language, def:desc-fo-semantics, def:P}
  For every first-order sentence $\varphi$, the language of the query
  $\mathcal{A} \mapsto (\mathcal{A} \models \varphi)$ is in $\cc{P}$.
\end{theorem}
\begin{proof}
  \uses{thm:desc-model-checking, thm:desc-encoding-roundtrip}
  Parse the input, then run the model checker. With $\varphi$ fixed, it makes
  $O(n^{\mathrm{qr}(\varphi)})$ recursive calls, each reading truth-table and
  constant bits at computable offsets. The same recursion with one
  logarithmic-size counter per quantifier would give the stronger bound
  $\cc{FO} \subseteq \cc{L}$.
\end{proof}

\section{First-order reductions}

\begin{definition}[First-order interpretation]
  \label{def:desc-interpretation}
  \lean{Complexity.DescriptiveComplexity.FOInterpretation, Complexity.DescriptiveComplexity.FOInterpretation.apply, Complexity.DescriptiveComplexity.FOInterpretation.idInterp, Complexity.DescriptiveComplexity.FOInterpretation.translateTerm, Complexity.DescriptiveComplexity.FOInterpretation.translate, Complexity.DescriptiveComplexity.FOInterpretation.comp}
  \leanok
  \uses{def:desc-fo-syntax, def:desc-structure, def:desc-substitution}
  A dimension-1 interpretation $I : V \to W$ consists of a $V$-formula with
  $\mathrm{ar}(R)$ free variables for each relation symbol $R$ of $W$, and a
  source constant for each constant of $W$. Applying $I$ to a $V$-structure gives
  a $W$-structure on the \emph{same} universe, in which $R$ holds exactly where
  its defining formula does. Formulas translate backwards along $I$, from $W$ to
  $V$: a target atom becomes its defining formula with the translated argument
  terms substituted, and a target constant becomes its source constant.
  Interpretations $U \to V$ and $V \to W$ compose to one $U \to W$ by translating
  the outer defining formulas. The identity interpretation defines each relation
  by its own atom and sends each constant to itself.
\end{definition}

\begin{theorem}[Transport theorem]
  \label{thm:desc-transport}
  \lean{Complexity.DescriptiveComplexity.FOInterpretation.translateTerm_eval, Complexity.DescriptiveComplexity.FOInterpretation.translate_sat, Complexity.DescriptiveComplexity.FOInterpretation.apply_comp, Complexity.DescriptiveComplexity.FOInterpretation.apply_idInterp}
  \leanok
  \uses{def:desc-interpretation, def:desc-fo-semantics}
  For every $V$-structure $\mathcal{A}$, environment $\sigma$, and $W$-formula
  $\varphi$, $I(\mathcal{A}), \sigma \models \varphi$ iff
  $\mathcal{A}, \sigma \models I^*\varphi$, and a translated term takes the same
  value in $\mathcal{A}$ as the original term in $I(\mathcal{A})$. Composition of
  interpretations acts as composition of the structure maps,
  $(I_2 \circ I_1)(\mathcal{A}) = I_2(I_1(\mathcal{A}))$, and the identity
  interpretation acts as the identity.
\end{theorem}
\begin{proof}
  \leanok
  \uses{thm:desc-substitution}
  Induct on $\varphi$. At a relation atom, substitute the translated argument
  terms into the defining formula. The composition law follows by applying the
  transport theorem to each defining formula.
\end{proof}

\begin{definition}[First-order reductions and projections]
  \label{def:desc-fo-reduction}
  \lean{Complexity.DescriptiveComplexity.FOReduces, Complexity.DescriptiveComplexity.FOInterpretation.IsQuantifierFree, Complexity.DescriptiveComplexity.FOProjReduces}
  \leanok
  \uses{def:desc-interpretation, def:desc-query}
  For queries $Q_1$ over $V$ and $Q_2$ over $W$, $Q_1 \le_{\cc{FO}} Q_2$ if some
  dimension-1 interpretation $I : V \to W$ has
  $Q_1(\mathcal{A}) \iff Q_2(I(\mathcal{A}))$ for all $V$-structures
  $\mathcal{A}$. An interpretation is \emph{quantifier-free} if every defining
  formula has quantifier rank $0$. $Q_1$ \emph{first-order projection} reduces to
  $Q_2$ if such an $I$ can be chosen quantifier-free. This condition is coarser
  than Immerman's projective form, in which each target bit depends on a single
  source bit.
\end{definition}

\begin{theorem}[Reductions are preorders]
  \label{thm:desc-fo-reduction-preorder}
  \lean{Complexity.DescriptiveComplexity.FOReduces.refl, Complexity.DescriptiveComplexity.FOReduces.trans, Complexity.DescriptiveComplexity.FOReduces.complement, Complexity.DescriptiveComplexity.FOProjReduces.refl, Complexity.DescriptiveComplexity.FOProjReduces.trans, Complexity.DescriptiveComplexity.FOProjReduces.toFOReduces, Complexity.DescriptiveComplexity.FOInterpretation.idInterp_isQuantifierFree, Complexity.DescriptiveComplexity.FOInterpretation.IsQuantifierFree.quantifierRank_translate, Complexity.DescriptiveComplexity.FOInterpretation.IsQuantifierFree.comp}
  \leanok
  \uses{def:desc-fo-reduction}
  First-order reducibility and first-order projection reducibility are reflexive
  and transitive, across vocabularies. Every projection is a reduction, and a
  reduction from $Q_1$ to $Q_2$ is, through the same interpretation, also a
  reduction from $\overline{Q_1}$ to $\overline{Q_2}$. The identity
  interpretation is quantifier-free, translation along a quantifier-free
  interpretation preserves the quantifier rank of every formula, and composites
  of quantifier-free interpretations are quantifier-free.
\end{theorem}
\begin{proof}
  \leanok
  \uses{thm:desc-transport, thm:desc-substitution}
  Use the identity interpretation and composition, with the composition law of
  the transport theorem. Translation along a quantifier-free interpretation
  replaces each atom by a quantifier-free formula with terms substituted, and
  substitution preserves quantifier rank, so quantifier rank is preserved and
  composites stay quantifier-free.
\end{proof}

\begin{theorem}[Definability is closed under reductions]
  \label{thm:desc-definable-reduction-closure}
  \lean{Complexity.DescriptiveComplexity.FODefinable.of_reduces, Complexity.DescriptiveComplexity.FODefinable.of_projReduces}
  \leanok
  \uses{def:desc-fo-reduction, def:desc-fo-definable}
  If $Q_1 \le_{\cc{FO}} Q_2$ and $Q_2$ is first-order definable, then so is $Q_1$.
  In particular, this holds when $Q_1$ first-order projection reduces to $Q_2$.
\end{theorem}
\begin{proof}
  \leanok
  \uses{thm:desc-transport, thm:desc-fo-reduction-preorder}
  Translate the sentence defining $Q_2$ backwards along the interpretation. A
  projection is a reduction.
\end{proof}

\begin{definition}[Dimension-$k$ interpretation]
  \label{def:desc-dim-k}
  \uses{def:desc-interpretation}
  A dimension-$k$ interpretation defines the target universe as a definable subset
  of $\mathrm{dom}^k$. Each target relation of arity $a$ is defined by a formula
  with $a k$ free variables. This requires a codec
  $\mathrm{Fin}(n^k) \simeq (\mathrm{Fin}(k) \to \mathrm{Fin}(n))$ and a way to
  assemble environments for the $a k$ variables.
\end{definition}

\begin{theorem}[Transport for dimension-$k$ interpretations]
  \label{thm:desc-dim-k-transport}
  \uses{def:desc-dim-k, def:desc-fo-semantics}
  \notready
  The transport theorem holds for dimension-$k$ interpretations. They compose
  (dimensions multiply), so dimension-$k$ reducibility is a preorder that
  contains dimension-1 reducibility.
\end{theorem}
\begin{proof}
  \uses{thm:desc-transport, thm:desc-substitution}
  Replace each target variable by a block of $k$ source variables and relativize
  quantifiers to the definable universe.
\end{proof}

\begin{theorem}[String-level reductions]
  \label{thm:desc-string-reduction}
  \uses{def:desc-fo-reduction, def:desc-query-language, def:reduction}
  If $Q_1 \le_{\cc{FO}} Q_2$, then the language of $Q_1$ polynomial-time
  many-one reduces to the language of $Q_2$.
\end{theorem}
\begin{proof}
  \uses{thm:desc-encoding-roundtrip, thm:desc-model-checking, thm:desc-transport, prop:desc-encoding-length}
  Decode the input. Compute each target truth table by model-checking the defining
  formulas, and re-encode. Send malformed inputs to a fixed non-encoding, such as
  the empty string: every encoding has length at least $n + 1 \ge 3$.
\end{proof}

\section{Capturing complexity classes}

\begin{definition}[Ordered first-order logic]
  \label{def:desc-ordered-fo}
  \uses{def:desc-fo-syntax, def:desc-structure}
  First-order logic extended with atoms for the canonical order $\le$, successor,
  minimum, and maximum of $\mathrm{Fin}(n)$. These atoms are interpreted by the
  existing meta-level helpers. The existing unordered syntax and its unrestricted
  invariance theorem stay unchanged. Ordered sentences are preserved only by maps
  that respect the canonical order, and they need not define order-independent
  queries.
\end{definition}

\begin{theorem}[Fagin, $\exists\cc{SO} \subseteq \cc{NP}$]
  \label{thm:desc-fagin-upper}
  \uses{def:desc-existential-so, def:desc-so-semantics, def:desc-so-definable, def:desc-query-language, def:NP}
  Let $\varphi$ be a second-order sentence in $\exists\cc{SO}$ form that defines
  $Q$. Then the language of $Q$ is in $\cc{NP}$.
\end{theorem}
\begin{proof}
  \uses{thm:desc-model-checking, thm:desc-fo-in-P, thm:desc-encoding-roundtrip}
  The certificate is the truth tables of the quantified relations, of polynomial
  total length $\sum_j n^{k_j}$. The verifier model-checks the FO matrix under
  the resulting relation environment, extending \texttt{evalB} to relation
  variables.
\end{proof}

\begin{theorem}[Fagin, $\cc{NP} \subseteq \exists\cc{SO}$]
  \label{thm:desc-fagin-lower}
  \uses{def:desc-existential-so, def:desc-so-semantics, def:desc-so-definable, def:desc-query-language, def:NP}
  Let $Q$ be order-independent with language in $\cc{NP}$. Then some second-order
  sentence in $\exists\cc{SO}$ form defines $Q$.
\end{theorem}
\begin{proof}
  \uses{thm:desc-so-definable, thm:desc-encoding-roundtrip}
  Existentially quantify a linear order together with a computation tableau of an
  $n^k$-time verifier, indexed by $k$-tuples. First-order-check the tableau's
  initial row (the encoding of the structure under the guessed order), its local
  transition windows, and acceptance. Order independence makes the guessed order
  harmless.
\end{proof}

\begin{theorem}[Fagin's theorem]
  \label{thm:desc-fagin}
  \uses{def:desc-existential-so, def:desc-query-language, def:NP}
  An order-independent query has its language in $\cc{NP}$ if and only if it is
  definable by an $\exists\cc{SO}$ sentence.
\end{theorem}
\begin{proof}
  \uses{thm:desc-fagin-upper, thm:desc-fagin-lower, thm:desc-so-definable}
  Combine the two inclusions. Order independence of $\exists\cc{SO}$-definable
  queries is already formalized.
\end{proof}

\begin{theorem}[$\cc{FO} \subseteq \cc{AC^0}$, nonuniform]
  \label{thm:desc-fo-ac0}
  \uses{def:desc-query-language, def:desc-fo-semantics, def:NC-AC, def:bitstring}
  For every first-order sentence $\varphi$, consider the characteristic family of
  the language of $\mathcal{A} \mapsto (\mathcal{A} \models \varphi)$: the Boolean
  function at input length $N$ is the indicator of the language restricted to
  $\{0,1\}^N$. This family lies in $\cc{AC^0}$.
\end{theorem}
\begin{proof}
  \uses{prop:desc-encoding-length, thm:desc-model-checking}
  The encoding length is strictly increasing in the universe size, so at most one
  size $n$ has encodings of length $N$. Build a circuit that checks the unary
  prefix and that each constant block is one-hot. Replace each quantifier by an
  unbounded fan-in gate of fan-in $n$, and each atom by an input bit, or by an
  unbounded fan-in selection over the one-hot block when an argument is a
  constant. The depth is $O(\lvert\varphi\rvert)$, a constant, and the size is
  polynomial in $n$, hence in $N$.
\end{proof}

\begin{theorem}[Barrington--Immerman--Straubing]
  \label{thm:desc-ac0-fo}
  \uses{def:desc-ordered-fo, def:NC-AC}
  \notready
  On ordered structures with a BIT predicate, the first-order definable queries
  are exactly those whose languages have $\mathsf{DLOGTIME}$-uniform $\cc{AC^0}$
  circuit families. This makes the previous theorem uniform and adds its
  converse.
\end{theorem}
\begin{proof}
  \uses{thm:desc-fo-ac0}
  One direction refines the circuit construction to a uniform one. For the other,
  describe the uniform circuit's gates by first-order formulas over tuples and
  evaluate the circuit by a first-order formula of depth proportional to the
  circuit's depth. Neither uniform $\cc{AC^0}$ nor BIT is defined yet.
\end{proof}

\begin{definition}[Least fixed-point logic]
  \label{def:desc-lfp}
  \uses{def:desc-fo-syntax, def:desc-fo-semantics}
  $\cc{FO}(\mathrm{LFP})$ extends first-order logic with
  $[\mathrm{lfp}_{R,\bar x}\, \varphi](\bar t)$ for $\varphi$ positive in $R$. It
  is interpreted as the least fixed point of the monotone operator that
  $\varphi$ induces on relations of the arity of $R$. That fixed point is reached
  after at most $n^{\mathrm{ar}(R)}$ iterations.
\end{definition}

\begin{theorem}[Immerman--Vardi]
  \label{thm:desc-immerman-vardi}
  \uses{def:desc-lfp, def:desc-ordered-fo, def:desc-query-language, def:P}
  \notready
  On ordered structures, a query is $\cc{FO}(\mathrm{LFP})$-definable if and only
  if its language is in $\cc{P}$.
\end{theorem}
\begin{proof}
  \uses{thm:desc-fo-in-P, thm:desc-encoding-roundtrip}
  $\subseteq$: iterate each fixed point at most polynomially many times. $\supseteq$:
  use the ordering to define the computation tableau of a polynomial-time machine
  by a least fixed point, as in the proof of Fagin's theorem but without guessing.
\end{proof}

\section{Open directions}

Immerman's exact projective reductions (each target bit determined by one source
bit under a quantifier-free guard) are not yet defined. Neither is a worked
completeness result under first-order projections. Both are prerequisites for
treating first-order projections as the reductions under which the standard
$\cc{AC^0}$-, $\cc{L}$-, and $\cc{NL}$-complete problems are complete.
